\section{Statistical Estimands and Uncertainty}
%% ============================================================

Let $Y_s^B$ and $Y_s^A$ denote final outcomes for paired training seed $s$. We estimate the mean channel-replacement effect as
\[
\hat{\Delta} = \frac{1}{10}\sum_{s=1}^{10}(Y_s^A - Y_s^B).
\]
Each bootstrap draw resamples ten paired seed indices with replacement. We report the 2.5th and 97.5th percentiles from \num{2000} draws using NumPy \texttt{default\_rng}, seed 20260920. Whole-seed resampling also gives pointwise learning-curve bands. Intervals condition on the posterior, simulator, hyperparameters, and evaluation protocol; they are not simultaneous bands.

For the counterfactual price-cut audits, we intervene on either firm ($0$ or $1$) at period $t_0\in\{64, 256\}$ using normalized action cuts of $0.1$ or $0.3$, giving eight probe types per seed. Because the action-to-price mapping has slope $0.75$, these normalized cuts correspond to unclipped price reductions of $0.075$ and $0.225$. Each probe uses 128 paired rollouts sharing the episode's posterior draw $\theta$ and Gaussian innovations $\varepsilon_t$; after the single forced cut at $t_0$, both agents return to their frozen deterministic policies through period $511$. We record both the signed sum of the rival impulse response $\Delta p_t^{\mathrm{rival}}$ after impact and the sum of the negative parts of its mean path (a descriptive downward-response area, not a structural measure of punishment). Similarly, the discounted continuation gain $\Delta\pi^{\mathrm{cont}}$ includes both firms' subsequent price adjustments and does not isolate the rival's causal contribution alone. In seed-level summaries, we average probe types within each seed before resampling seeds: cut-specific means average four probe types per seed, while overall rival-response area averages all eight. Seeds receive equal weight; the 128 rollout repetitions and four within-cut probe types are not independent training replications.

\begin{table}[t]
\centering
\caption{Seed-bootstrap intervals for cut-specific full gain and its continuation component. Intervals are pointwise and unadjusted.}
\label{tab:bootstrap_ci}
\small
\sisetup{table-format=-1.6}
\begin{tabular}{@{}cc l SSS@{}}
\toprule
Arm & Cut & Quantity & {Mean} & {Lower 95\%} & {Upper 95\%} \\
\midrule
B & 0.1 & Full gain      & -0.005192 & -0.005855 & -0.004542 \\
B & 0.1 & Continuation   & -0.000005 & -0.000010 & -0.000001 \\
B & 0.3 & Full gain      & -0.051322 & -0.053372 & -0.049286 \\
B & 0.3 & Continuation   & -0.000005 & -0.000010 & -0.000001 \\
A & 0.1 & Full gain      & -0.004800 & -0.005393 & -0.004343 \\
A & 0.1 & Continuation   &  0.000000 &  0.000000 &  0.000000 \\
A & 0.3 & Full gain      & -0.050273 & -0.052023 & -0.048937 \\
A & 0.3 & Continuation   &  0.000000 &  0.000000 &  0.000000 \\
\bottomrule
\end{tabular}
\end{table}

\subsection{Conditional Intervals in the Complete Probe Tables}

Each individual probe file also records a \num{1000}-resample bootstrap interval for the mean across its 128 paired rollouts. The bootstrap generator is initialized with the audit seed plus \num{54321}. Tables~\ref{tab:probe_B01}--\ref{tab:probe_A03} report those stored mean intervals. Their interpretation is conditional on one trained policy pair. The separately stored rollout quantiles named \texttt{deviation\_gain\_q025} and \texttt{deviation\_gain\_q975} are not confidence limits for a mean and are not used in those tables.

In Tables~\ref{tab:probe_B01}--\ref{tab:probe_A03}, seed suffixes 01--10 identify 20261001--20261010; $t$ is intervention time. Immediate denotes impact gain; $G$ includes the remaining episode. Area sums the rival's mean price response after impact, in price-period units. Zeros are rounded except exact action-independent responses. Lower and upper limits are unadjusted conditional mean intervals.

Within-seed probes share policies and random-stream construction. Fixed diagnostic streams across checkpoints likewise do not create independent observations. The 100 checkpoints cannot be treated as 100 independent training samples.

\subsection{What an Unresolved Difference Means}

The price-elevation interval includes zero and effects of either sign. Without a tested equivalence margin, this does not establish equivalence. The four positive conditional small-cut means are descriptive, not multiplicity-adjusted discoveries. Reporting both averages and individual probes preserves this heterogeneity.

%% ============================================================
